\section{Introduction}
Over the last couple of years, vehicle monitoring has been a necessary measure to make vehicles and drivers safe, cut down on emissions to tackle climate change , and provide optimal performance. The significance of protecting human life and the necessity to obtain genuine information at the time of accidents for forensic analysis reinforce the importance of efficient monitoring systems. To serve this purpose, we created a vehicle monitoring logger on the basis of On-Board Diagnostics (OBD-II) port. The OBD-II port is a 16-pin standardized connector found in every contemporary vehicle,developed in the United States Of America in 1996, by the  Society of Automotive Engineers (SAE). It was imposed to enable regulation  of vehicle emission so as to ensure Environmental Protection Agency(EPA) standards are met. It is a diagnostic port used to monitor real-time readings of vehicle sensors and the Diagnostic Trouble Codes (DTCs) from the Electronic Control Unit (ECU), making it a valuable resource for monitoring vehicle performance.Ecu is a the main controller of the vehicle where all controls and actions are done.

OBD II port supports communication protocols such as ISO15765-4 (CAN-BUS), ISO14230-4 (KWP2000), ISO 9141-2, SAE J1850 VPW, and SAE J1850 PWM.  This study is centered on the ISO 15765-4 (CAN-BUS) protocol that has been required on vehicles manufactured from 1996.
The message request is send in a hexadecimal to the protocols and it gets interpreted according  to one of the five OBD II protocols and sends back a hexadecimal response.Obd II uses two types of codes to request ECU data that are Diagnostic Trouble Code (DTCs) and Parameter Identifiers (PIDs).DTC are used to diagnostic malfunctions in various subsystem of the vehicle and PIDs are used to measure real time parameters 
